\documentclass[10pt]{article}

% =========================
% Language and encoding
% =========================
\usepackage[english]{babel}
\usepackage[utf8]{inputenc}
\usepackage[T1]{fontenc}
\usepackage{lmodern}

% =========================
% Page layout
% =========================
\usepackage[landscape,margin=0.9cm]{geometry}
\usepackage{multicol}
\setlength{\columnsep}{0.75cm}
\setlength{\columnseprule}{0.35pt}
\setlength{\parindent}{0pt}
\setlength{\parskip}{3pt}
\raggedcolumns

% =========================
% Math, color, links, lists
% =========================
\usepackage{amsmath}
\usepackage{amssymb}
\usepackage{xcolor}
\usepackage{hyperref}
\usepackage{enumitem}
\usepackage{titlesec}
\usepackage{listings}
\usepackage{needspace}

\hypersetup{
    colorlinks=true,
    linkcolor=blue!45!black,
    urlcolor=blue!45!black
}

% =========================
% Section style
% =========================
\titleformat{\section}{\Large\bfseries\color{blue!45!black}}{\thesection}{0.7em}{}
\titleformat{\subsection}{\normalsize\bfseries\color{black}}{\thesubsection}{0.6em}{}
\titlespacing*{\section}{0pt}{7pt}{4pt}
\titlespacing*{\subsection}{0pt}{4pt}{2pt}

\setlist[itemize]{leftmargin=*, itemsep=1pt, topsep=2pt}

% =========================
% Code style
% =========================
\lstdefinestyle{pystyle}{
    language=Python,
    basicstyle=\ttfamily\scriptsize,
    keywordstyle=\bfseries\color{blue!55!black},
    commentstyle=\itshape\color{gray!70!black},
    stringstyle=\color{green!35!black},
    breaklines=true,
    breakatwhitespace=false,
    columns=fullflexible,
    keepspaces=true,
    showstringspaces=false,
    tabsize=2,
    frame=single,
    rulecolor=\color{black!20},
    framerule=0.3pt,
    xleftmargin=1pt,
    xrightmargin=1pt,
    aboveskip=3pt,
    belowskip=3pt
}
\lstset{style=pystyle}

\newcommand{\code}[1]{\texttt{#1}}
\newenvironment{cpblock}[1]{\par\Needspace{18\baselineskip}\begin{minipage}{\linewidth}\subsection{#1}}{\end{minipage}\par\vspace{5pt}}

\begin{document}
\pagenumbering{gobble}

% =========================
% Cover
% =========================
\begin{titlepage}
    \centering
    \vspace*{2cm}
    {\Huge\bfseries Python Competitive Programming Reference Book\par}
    \vspace{0.7cm}
    {\Huge Arteaga Adán\par}
    \vspace{1.2cm}
    {\Huge Algorithms Sheet - Python Edition\par}
    \vspace{0.5cm}
    {\Large Clean, contest-ready, interview-friendly snippets\par}
    \vfill
    {\large \textbf{Python version}\par}
    
    \vfill
    {\large Python version for competitive programming\par}
\end{titlepage}

\tableofcontents
\newpage
\pagenumbering{arabic}

% =========================
\section{Python Contest Setup}
% =========================
\begin{cpblock}{Fast Input}
Use this when input is large. It is usually faster than repeated \code{input()}.
\begin{lstlisting}
import sys
input = sys.stdin.readline

n = int(input())
a = list(map(int, input().split()))
\end{lstlisting}
\end{cpblock}

\begin{cpblock}{Read All Tokens}
Best for very large input with many integers.
\begin{lstlisting}
import sys

data = sys.stdin.buffer.read().split()
it = iter(data)

n = int(next(it))
a = [int(next(it)) for _ in range(n)]
\end{lstlisting}
\end{cpblock}

\begin{cpblock}{Useful Imports}
\begin{lstlisting}
import sys
from math import gcd, isqrt, inf
from collections import deque, defaultdict, Counter
from bisect import bisect_left, bisect_right
from heapq import heappush, heappop
from itertools import accumulate, permutations, combinations
\end{lstlisting}
\end{cpblock}

\begin{cpblock}{Recursion Limit}
For deep DFS, increase the recursion limit. If recursion still risks stack issues, use iterative DFS.
\begin{lstlisting}
import sys
sys.setrecursionlimit(1_000_000)
\end{lstlisting}
\end{cpblock}

\begin{cpblock}{Infinity and Constants}
\begin{lstlisting}
INF = 10**30
MOD = 998244353
# Common alternative:
MOD2 = 10**9 + 7
\end{lstlisting}
\end{cpblock}

\begin{cpblock}{Python Contest Notes}
\begin{itemize}
    \item Prefer \code{sys.stdin.readline} or \code{buffer.read().split()} for heavy input.
    \item Use \code{deque} for queues, not \code{list.pop(0)}.
    \item Use \code{heapq} for priority queues.
    \item Use \code{bisect} for binary search on sorted arrays.
    \item Use \code{dict}, \code{set}, and \code{Counter} heavily.
    \item Prefer iterative DFS for very deep graphs.
\end{itemize}
\end{cpblock}

\newpage

% =========================
\section{Complexities and Python Structures}
% =========================
\begin{cpblock}{Common Time Complexities}
\begin{itemize}
    \item $O(1)$ - constant: index access, hash lookup average case.
    \item $O(\log n)$ - logarithmic: binary search, heap push/pop.
    \item $O(\sqrt n)$ - square root: divisors, simple primality checks.
    \item $O(n)$ - linear: one pass over an array.
    \item $O(n \log n)$ - sorting, many greedy solutions.
    \item $O(n^2)$ - nested loops, some dynamic programming.
    \item $O(2^n)$ - subsets and bitmask DP.
    \item $O(n!)$ - permutations.
\end{itemize}
\end{cpblock}

\begin{cpblock}{Python Data Structures}
\begin{itemize}
    \item \code{list}: append/pop back $O(1)$ amortized, insert/delete middle $O(n)$.
    \item \code{dict/set}: insert/find average $O(1)$, worst $O(n)$.
    \item \code{deque}: push/pop both ends $O(1)$.
    \item \code{heapq}: push/pop $O(\log n)$, min element $O(1)$.
    \item \code{bisect}: search in sorted list $O(\log n)$, insertion still $O(n)$.
    \item \code{Counter}: frequency map.
\end{itemize}
\end{cpblock}

\begin{cpblock}{Standard Constraints}
\begin{itemize}
    \item Python is slower than compiled languages, so constants matter.
    \item $n=10^5$: usually $O(n)$ or $O(n\log n)$.
    \item $n=500$: $O(n^3)$ may pass depending on constants.
    \item $n=20$: $O(2^n \cdot n)$ is common.
    \item $10^7$ simple operations can be risky but possible with optimized code.
\end{itemize}
\end{cpblock}

\begin{cpblock}{Important Algorithms}
\begin{itemize}
    \item GCD: $O(\log(\min(a,b)))$
    \item Binary exponentiation: $O(\log exp)$
    \item Sieve of Eratosthenes: $O(n\log\log n)$
    \item Prime factorization by trial division: $O(\sqrt n)$
    \item DFS/BFS: $O(V+E)$
    \item Dijkstra: $O(E\log V)$
    \item Floyd-Warshall: $O(V^3)$
    \item DSU: $O(\alpha(n))$ amortized
    \item Segment Tree query/update: $O(\log n)$
    \item Fenwick Tree query/update: $O(\log n)$
    \item KMP, Z-function, Manacher: $O(n)$
\end{itemize}
\end{cpblock}

\newpage

% =========================
\section{Input, Strings, and Parsing Basics}
% =========================
\begin{cpblock}{Read a Full Line}
Equivalent to reading a whole line including spaces.
\begin{lstlisting}
s = input().rstrip('\n')
\end{lstlisting}
\end{cpblock}

\begin{cpblock}{Split a Line}
Equivalent to tokenizing a string by whitespace.
\begin{lstlisting}
s = "10 20 30"
a, b, c = map(int, s.split())
\end{lstlisting}
\end{cpblock}

\begin{cpblock}{Character and Word Tools}
\begin{lstlisting}
from collections import Counter

s = "abacaba"
freq = Counter(s)

words = "one two three".split()
joined = "-".join(words)
\end{lstlisting}
\end{cpblock}

\begin{cpblock}{Sorting Strings}
\begin{lstlisting}
s = "dcba"
t = ''.join(sorted(s))

arr = ["aaa", "b", "cc"]
arr.sort(key=len)
\end{lstlisting}
\end{cpblock}

\newpage

% =========================
\section{Binary Search}
% =========================
\begin{cpblock}{Basic Binary Search}
\begin{lstlisting}
l, r = 0, n - 1
ans = -1

while l <= r:
    mid = (l + r) // 2
    if a[mid] == x:
        ans = mid
        break
    if a[mid] < x:
        l = mid + 1
    else:
        r = mid - 1
\end{lstlisting}
\end{cpblock}

\begin{cpblock}{Binary Search on Answer}
Find the minimum value satisfying a monotonic predicate.
\begin{lstlisting}
l, r = 0, 10**18
ans = r

while l <= r:
    mid = (l + r) // 2
    if check(mid):
        ans = mid
        r = mid - 1
    else:
        l = mid + 1
\end{lstlisting}
\end{cpblock}

\begin{cpblock}{bisect Left and Right}
\begin{lstlisting}
from bisect import bisect_left, bisect_right

idx = bisect_left(a, x)   # first position with a[i] >= x
idx2 = bisect_right(a, x) # first position with a[i] > x

count_x = bisect_right(a, x) - bisect_left(a, x)
exists = bisect_left(a, x) < len(a) and a[bisect_left(a, x)] == x
\end{lstlisting}
\end{cpblock}

\begin{cpblock}{Last True / First True Pattern}
\begin{lstlisting}
# first True over [lo, hi)
lo, hi = 0, n
while lo < hi:
    mid = (lo + hi) // 2
    if ok(mid):
        hi = mid
    else:
        lo = mid + 1
ans = lo
\end{lstlisting}
\end{cpblock}

\newpage

% =========================
\section{Math and Number Theory}
% =========================
\begin{cpblock}{GCD and LCM}
\begin{lstlisting}
from math import gcd

def lcm(a, b):
    if a == 0 or b == 0:
        return 0
    return abs(a // gcd(a, b) * b)
\end{lstlisting}
\textbf{Properties:}
\begin{itemize}
    \item $gcd(a,b) \cdot lcm(a,b)=|a\cdot b|$
    \item $gcd(a,b)=gcd(b,a \bmod b)$
    \item $gcd(a,b,c)=gcd(a,gcd(b,c))$
\end{itemize}
\end{cpblock}

\begin{cpblock}{Manual Euclidean GCD}
\begin{lstlisting}
def gcd_manual(a, b):
    while b:
        a, b = b, a % b
    return abs(a)
\end{lstlisting}
\end{cpblock}

\begin{cpblock}{Classic Sieve + Factorization}
\begin{lstlisting}
def sieve(n):
    is_prime = [True] * (n + 1)
    is_prime[0:2] = [False, False]
    primes = []
    for i in range(2, n + 1):
        if is_prime[i]:
            primes.append(i)
            if i * i <= n:
                for j in range(i * i, n + 1, i):
                    is_prime[j] = False
    return is_prime, primes

def factorize(x, primes):
    f = []
    for p in primes:
        if p * p > x:
            break
        while x % p == 0:
            f.append(p)
            x //= p
    if x > 1:
        f.append(x)
    return f
\end{lstlisting}
\end{cpblock}

\begin{cpblock}{Linear SPF Sieve + Factorization}
\begin{lstlisting}
def linear_sieve_spf(n):
    spf = [0] * (n + 1)
    primes = []
    for i in range(2, n + 1):
        if spf[i] == 0:
            spf[i] = i
            primes.append(i)
        for p in primes:
            if p > spf[i] or i * p > n:
                break
            spf[i * p] = p
    return spf, primes

def factorize_spf(x, spf):
    f = []
    while x > 1:
        f.append(spf[x])
        x //= spf[x]
    return f
\end{lstlisting}
\end{cpblock}

\begin{cpblock}{Euler's Totient Function}
\begin{lstlisting}
def phi_single(n):
    res = n
    p = 2
    while p * p <= n:
        if n % p == 0:
            while n % p == 0:
                n //= p
            res -= res // p
        p += 1
    if n > 1:
        res -= res // n
    return res
\end{lstlisting}
\end{cpblock}

\begin{cpblock}{Totient Sieve}
\begin{lstlisting}
def phi_sieve(n):
    phi = list(range(n + 1))
    for i in range(2, n + 1):
        if phi[i] == i:  # prime
            for j in range(i, n + 1, i):
                phi[j] -= phi[j] // i
    return phi
\end{lstlisting}
\end{cpblock}

\begin{cpblock}{Notable Sums}
\[
\sum_{i=1}^{n} i = \frac{n(n+1)}{2}
\]
\[
\sum_{i=1}^{n} i^2 = \frac{n(n+1)(2n+1)}{6}
\]
\begin{lstlisting}
sum_1_to_n = n * (n + 1) // 2
sum_squares = n * (n + 1) * (2 * n + 1) // 6
\end{lstlisting}
\end{cpblock}

\newpage

% =========================
\section{Modular Arithmetic and Combinatorics}
% =========================
\begin{cpblock}{Binary Exponentiation}
Python already has optimized modular exponentiation through \code{pow(a,b,mod)}.
\begin{lstlisting}
def binpow(a, b, mod):
    res = 1
    a %= mod
    while b:
        if b & 1:
            res = res * a % mod
        a = a * a % mod
        b >>= 1
    return res

ans = pow(a, b, MOD)
\end{lstlisting}
\end{cpblock}

\begin{cpblock}{Modular Inverse}
Works when \code{MOD} is prime and \code{a} is not divisible by \code{MOD}.
\begin{lstlisting}
def mod_inverse(a, mod):
    return pow(a, mod - 2, mod)
\end{lstlisting}
\end{cpblock}

\begin{cpblock}{Modulo Properties}
\[
(a+b)\bmod m=((a\bmod m)+(b\bmod m))\bmod m
\]
\[
(a-b)\bmod m=((a\bmod m)-(b\bmod m))\bmod m
\]
\[
(a\cdot b)\bmod m=((a\bmod m)(b\bmod m))\bmod m
\]
In Python, \code{x \% MOD} already returns a non-negative result for positive \code{MOD}.
\end{cpblock}

\begin{cpblock}{Modular Factorials}
\begin{lstlisting}
MOD = 998244353
MAXN = 10**6 + 5
fact = [1] * MAXN

for i in range(1, MAXN):
    fact[i] = fact[i - 1] * i % MOD
\end{lstlisting}
\end{cpblock}

\begin{cpblock}{Modular Combinations}
\begin{lstlisting}
def nCr(n, r):
    if r < 0 or r > n:
        return 0
    return fact[n] * pow(fact[r], MOD - 2, MOD) % MOD * pow(fact[n - r], MOD - 2, MOD) % MOD
\end{lstlisting}
\[
C(n,r)=\frac{n!}{r!(n-r)!}
\qquad
P(n,r)=\frac{n!}{(n-r)!}
\]
\end{cpblock}

\begin{cpblock}{Precomputed Inverse Factorials}
Faster when there are many queries.
\begin{lstlisting}
fact = [1] * (MAXN + 1)
invfact = [1] * (MAXN + 1)

for i in range(1, MAXN + 1):
    fact[i] = fact[i - 1] * i % MOD

invfact[MAXN] = pow(fact[MAXN], MOD - 2, MOD)
for i in range(MAXN, 0, -1):
    invfact[i - 1] = invfact[i] * i % MOD

def comb(n, r):
    if r < 0 or r > n:
        return 0
    return fact[n] * invfact[r] % MOD * invfact[n - r] % MOD
\end{lstlisting}
\end{cpblock}

\newpage

% =========================
\section{Prefix Sums and Difference Arrays}
% =========================
\begin{cpblock}{Prefix Sum 1D}
\begin{lstlisting}
pref = [0] * (n + 1)
for i in range(n):
    pref[i + 1] = pref[i] + a[i]

def range_sum(l, r):  # inclusive, 0-indexed
    return pref[r + 1] - pref[l]
\end{lstlisting}
\end{cpblock}

\begin{cpblock}{Prefix XOR}
\begin{lstlisting}
px = [0] * (n + 1)
for i in range(n):
    px[i + 1] = px[i] ^ a[i]

def range_xor(l, r):
    return px[r + 1] ^ px[l]
\end{lstlisting}
\end{cpblock}

\begin{cpblock}{Prefix Sum 2D}
\begin{lstlisting}
p = [[0] * (m + 1) for _ in range(n + 1)]

for i in range(1, n + 1):
    for j in range(1, m + 1):
        p[i][j] = a[i - 1][j - 1] + p[i - 1][j] + p[i][j - 1] - p[i - 1][j - 1]

def rect_sum(x1, y1, x2, y2):
    # inclusive, 0-indexed
    x1 += 1; y1 += 1; x2 += 1; y2 += 1
    return p[x2][y2] - p[x1 - 1][y2] - p[x2][y1 - 1] + p[x1 - 1][y1 - 1]
\end{lstlisting}
\end{cpblock}

\begin{cpblock}{Difference Array}
Range update $[l,r]$ with value $v$.
\begin{lstlisting}
d = [0] * (n + 1)

def add_range(l, r, v):
    d[l] += v
    if r + 1 < len(d):
        d[r + 1] -= v

cur = 0
for i in range(n):
    cur += d[i]
    a[i] += cur
\end{lstlisting}
\end{cpblock}

\newpage

% =========================
\section{Bit Manipulation}
% =========================
\begin{cpblock}{Basic Operations}
\begin{lstlisting}
mask | (1 << i)       # set bit i
mask & (1 << i)       # check bit i
mask ^ (1 << i)       # toggle bit i
mask & ~(1 << i)      # clear bit i
\end{lstlisting}
\end{cpblock}

\begin{cpblock}{Least Significant Bit}
\begin{lstlisting}
mask & -mask          # isolate LSB
mask & (mask - 1)     # remove LSB
\end{lstlisting}
\end{cpblock}

\begin{cpblock}{Python Bit Built-ins}
\begin{lstlisting}
x.bit_count()         # number of 1 bits
x.bit_length()        # bits needed to represent x
x.bit_length() - 1    # floor(log2(x)), if x > 0
\end{lstlisting}
\end{cpblock}

\begin{cpblock}{All Bits Set of Size n}
\begin{lstlisting}
full = (1 << n) - 1
\end{lstlisting}
\end{cpblock}

\begin{cpblock}{Check Power of Two}
\begin{lstlisting}
def is_power_of_two(n):
    return n > 0 and (n & (n - 1)) == 0
\end{lstlisting}
\end{cpblock}

\begin{cpblock}{XOR Properties}
\begin{itemize}
    \item $a \oplus a = 0$
    \item $a \oplus 0 = a$
    \item If $a \oplus b = c$, then $a \oplus c = b$.
\end{itemize}
\end{cpblock}

\begin{cpblock}{Iterating All Subsets}
\begin{lstlisting}
for mask in range(1 << n):
    for j in range(n):
        if mask & (1 << j):
            pass  # j is in subset
\end{lstlisting}
\end{cpblock}

\begin{cpblock}{Iterating Submasks}
\begin{lstlisting}
sub = mask
while sub:
    # sub is a submask of mask
    sub = (sub - 1) & mask
# excludes empty submask 0
\end{lstlisting}
\end{cpblock}

\begin{cpblock}{Subsets of Size K - Gosper's Hack}
\begin{lstlisting}
set_mask = (1 << k) - 1
while set_mask < (1 << n):
    # process set_mask
    c = set_mask & -set_mask
    r = set_mask + c
    set_mask = (((r ^ set_mask) >> 2) // c) | r
\end{lstlisting}
\end{cpblock}

\newpage

% =========================
\section{Core Graph Algorithms}
% =========================
\begin{cpblock}{Adjacency List}
\begin{lstlisting}
adj = [[] for _ in range(n)]

# undirected edge
adj[u].append(v)
adj[v].append(u)

# directed edge
adj[u].append(v)
\end{lstlisting}
\end{cpblock}

\begin{cpblock}{DFS Recursive}
\begin{lstlisting}
import sys
sys.setrecursionlimit(1_000_000)

vis = [False] * n

def dfs(u):
    vis[u] = True
    for v in adj[u]:
        if not vis[v]:
            dfs(v)
\end{lstlisting}
\end{cpblock}

\begin{cpblock}{DFS Iterative}
Safer when recursion depth is dangerous.
\begin{lstlisting}
def dfs_iter(start):
    vis = [False] * n
    stack = [start]
    vis[start] = True
    while stack:
        u = stack.pop()
        for v in adj[u]:
            if not vis[v]:
                vis[v] = True
                stack.append(v)
    return vis
\end{lstlisting}
\end{cpblock}

\begin{cpblock}{BFS}
\begin{lstlisting}
from collections import deque

def bfs(s):
    vis = [False] * n
    q = deque([s])
    vis[s] = True
    while q:
        u = q.popleft()
        for v in adj[u]:
            if not vis[v]:
                vis[v] = True
                q.append(v)
    return vis
\end{lstlisting}
\end{cpblock}

\begin{cpblock}{BFS Distances}
\begin{lstlisting}
from collections import deque

def bfs_dist(s):
    dist = [-1] * n
    q = deque([s])
    dist[s] = 0
    while q:
        u = q.popleft()
        for v in adj[u]:
            if dist[v] == -1:
                dist[v] = dist[u] + 1
                q.append(v)
    return dist
\end{lstlisting}
\end{cpblock}

\begin{cpblock}{Grid Directions}
\begin{lstlisting}
dx = [-1, 1, 0, 0]
dy = [0, 0, -1, 1]

for k in range(4):
    nx = x + dx[k]
    ny = y + dy[k]
    if 0 <= nx < R and 0 <= ny < C:
        pass
\end{lstlisting}
\end{cpblock}

\begin{cpblock}{Graph Properties}
\begin{itemize}
    \item Tree edges: $n-1$.
    \item Complete graph edges: $\frac{n(n-1)}{2}$.
    \item Handshaking lemma: $\sum deg(v)=2|E|$.
\end{itemize}
\end{cpblock}

\newpage

% =========================
\section{Graph Coloring and Cycle Detection}
% =========================
\begin{cpblock}{Bipartite Coloring}
2-color an undirected graph. If impossible, the graph is not bipartite.
\begin{lstlisting}
from collections import deque

def bipartite_color(adj):
    n = len(adj)
    color = [-1] * n
    for s in range(n):
        if color[s] != -1:
            continue
        color[s] = 0
        q = deque([s])
        while q:
            u = q.popleft()
            for v in adj[u]:
                if color[v] == -1:
                    color[v] = color[u] ^ 1
                    q.append(v)
                elif color[v] == color[u]:
                    return False, color
    return True, color
\end{lstlisting}
\end{cpblock}

\begin{cpblock}{Greedy Graph Coloring}
Assigns a valid coloring, not necessarily minimum.
\begin{lstlisting}
def greedy_coloring(adj):
    n = len(adj)
    color = [-1] * n
    for u in range(n):
        used = {color[v] for v in adj[u] if color[v] != -1}
        c = 0
        while c in used:
            c += 1
        color[u] = c
    return color
\end{lstlisting}
\end{cpblock}

\begin{cpblock}{Cycle Detection - Undirected Graph}
\begin{lstlisting}
def has_cycle_undirected(adj):
    n = len(adj)
    vis = [False] * n

    def dfs(u, parent):
        vis[u] = True
        for v in adj[u]:
            if not vis[v]:
                if dfs(v, u):
                    return True
            elif v != parent:
                return True
        return False

    for i in range(n):
        if not vis[i] and dfs(i, -1):
            return True
    return False
\end{lstlisting}
\end{cpblock}

\begin{cpblock}{Cycle Detection - Directed Graph}
Use states: 0 = unvisited, 1 = visiting, 2 = done.
\begin{lstlisting}
def has_cycle_directed(adj):
    n = len(adj)
    state = [0] * n

    def dfs(u):
        state[u] = 1
        for v in adj[u]:
            if state[v] == 1:
                return True
            if state[v] == 0 and dfs(v):
                return True
        state[u] = 2
        return False

    for i in range(n):
        if state[i] == 0 and dfs(i):
            return True
    return False
\end{lstlisting}
\end{cpblock}

\newpage

% =========================
\section{Shortest Paths and DSU}
% =========================
\begin{cpblock}{Dijkstra}
Graph format: \code{adj[u] = [(v, w), ...]}.
\begin{lstlisting}
from heapq import heappush, heappop

INF = 10**30

def dijkstra(adj, s):
    n = len(adj)
    dist = [INF] * n
    dist[s] = 0
    pq = [(0, s)]

    while pq:
        d, u = heappop(pq)
        if d != dist[u]:
            continue
        for v, w in adj[u]:
            nd = d + w
            if nd < dist[v]:
                dist[v] = nd
                heappush(pq, (nd, v))
    return dist
\end{lstlisting}
\end{cpblock}

\begin{cpblock}{Disjoint Set Union}
\begin{lstlisting}
class DSU:
    def __init__(self, n):
        self.parent = list(range(n))
        self.sz = [1] * n

    def find(self, x):
        while self.parent[x] != x:
            self.parent[x] = self.parent[self.parent[x]]
            x = self.parent[x]
        return x

    def unite(self, a, b):
        ra, rb = self.find(a), self.find(b)
        if ra == rb:
            return False
        if self.sz[ra] < self.sz[rb]:
            ra, rb = rb, ra
        self.parent[rb] = ra
        self.sz[ra] += self.sz[rb]
        return True
\end{lstlisting}
\end{cpblock}

\begin{cpblock}{Floyd-Warshall}
All-pairs shortest paths. Use only for small $n$.
\begin{lstlisting}
for k in range(n):
    for i in range(n):
        dik = dist[i][k]
        for j in range(n):
            nd = dik + dist[k][j]
            if nd < dist[i][j]:
                dist[i][j] = nd
\end{lstlisting}
\end{cpblock}

\newpage

% =========================
\section{String Matching Algorithms}
% =========================
\begin{cpblock}{KMP Prefix Function}
\begin{lstlisting}
def prefix_function(s):
    n = len(s)
    pi = [0] * n
    for i in range(1, n):
        j = pi[i - 1]
        while j > 0 and s[i] != s[j]:
            j = pi[j - 1]
        if s[i] == s[j]:
            j += 1
        pi[i] = j
    return pi
\end{lstlisting}
\end{cpblock}

\begin{cpblock}{KMP Pattern Matching}
Find all occurrences of pattern \code{p} in text \code{t}.
\begin{lstlisting}
def kmp_search(t, p):
    s = p + '#' + t
    pi = prefix_function(s)
    ans = []
    m = len(p)
    for i in range(m + 1, len(s)):
        if pi[i] == m:
            ans.append(i - 2 * m)
    return ans
\end{lstlisting}
\end{cpblock}

\begin{cpblock}{Z-Function}
\begin{lstlisting}
def z_function(s):
    n = len(s)
    z = [0] * n
    l = r = 0
    for i in range(1, n):
        if i <= r:
            z[i] = min(r - i + 1, z[i - l])
        while i + z[i] < n and s[z[i]] == s[i + z[i]]:
            z[i] += 1
        if i + z[i] - 1 > r:
            l, r = i, i + z[i] - 1
    return z
\end{lstlisting}
\end{cpblock}

\begin{cpblock}{Z-Function Pattern Matching}
\begin{lstlisting}
def z_search(t, p):
    s = p + '#' + t
    z = z_function(s)
    m = len(p)
    return [i - m - 1 for i in range(m + 1, len(s)) if z[i] >= m]
\end{lstlisting}
\end{cpblock}

\begin{cpblock}{Manacher - Odd Palindromes}
\code{d1[i]} = radius of odd palindrome centered at \code{i}.
\begin{lstlisting}
def manacher_odd(s):
    n = len(s)
    d1 = [0] * n
    l, r = 0, -1
    for i in range(n):
        k = 1 if i > r else min(d1[l + r - i], r - i + 1)
        while 0 <= i - k and i + k < n and s[i - k] == s[i + k]:
            k += 1
        d1[i] = k
        if i + k - 1 > r:
            l, r = i - k + 1, i + k - 1
    return d1
\end{lstlisting}
\end{cpblock}

\begin{cpblock}{Manacher - Even Palindromes}
\code{d2[i]} = radius of even palindrome centered between \code{i-1} and \code{i}.
\begin{lstlisting}
def manacher_even(s):
    n = len(s)
    d2 = [0] * n
    l, r = 0, -1
    for i in range(n):
        k = 0 if i > r else min(d2[l + r - i + 1], r - i + 1)
        while 0 <= i - k - 1 and i + k < n and s[i - k - 1] == s[i + k]:
            k += 1
        d2[i] = k
        if i + k - 1 > r:
            l, r = i - k, i + k - 1
    return d2
\end{lstlisting}
\end{cpblock}

\newpage

% =========================
\section{Dynamic Programming and Sequences}
% =========================
\begin{cpblock}{LIS - Longest Increasing Subsequence}
$O(n\log n)$ length only.
\begin{lstlisting}
from bisect import bisect_left, bisect_right

def lis_length(a):
    lis = []
    for x in a:
        i = bisect_left(lis, x)
        if i == len(lis):
            lis.append(x)
        else:
            lis[i] = x
    return len(lis)

# For longest non-decreasing subsequence, use bisect_right.
\end{lstlisting}
\end{cpblock}

\begin{cpblock}{Maximum Subarray Sum - Kadane}
Works even if all numbers are negative.
\begin{lstlisting}
def kadane(a):
    best = -10**30
    cur = 0
    for x in a:
        cur = max(x, cur + x)
        best = max(best, cur)
    return best
\end{lstlisting}
\end{cpblock}

\begin{cpblock}{Basic 1D DP Pattern}
\begin{lstlisting}
dp = [0] * (n + 1)
dp[0] = 1

for i in range(1, n + 1):
    # transition here
    dp[i] = dp[i - 1]
\end{lstlisting}
\end{cpblock}

\begin{cpblock}{Bitmask DP Skeleton}
\begin{lstlisting}
INF = 10**30
dp = [INF] * (1 << n)
dp[0] = 0

for mask in range(1 << n):
    for i in range(n):
        if not (mask & (1 << i)):
            new_mask = mask | (1 << i)
            dp[new_mask] = min(dp[new_mask], dp[mask] + cost[i])
\end{lstlisting}
\end{cpblock}

\newpage

% =========================
\section{Advanced Data Structures}
% =========================
\begin{cpblock}{Fenwick Tree - BIT}
1-indexed structure for prefix sums.
\begin{lstlisting}
class Fenwick:
    def __init__(self, n):
        self.n = n
        self.bit = [0] * (n + 1)

    def add(self, idx, val):
        while idx <= self.n:
            self.bit[idx] += val
            idx += idx & -idx

    def sum(self, idx):
        res = 0
        while idx > 0:
            res += self.bit[idx]
            idx -= idx & -idx
        return res

    def range_sum(self, l, r):
        return self.sum(r) - self.sum(l - 1)
\end{lstlisting}
\end{cpblock}

\begin{cpblock}{Segment Tree - Range Sum}
Iterative version: compact and fast in Python.
\begin{lstlisting}
class SegTree:
    def __init__(self, a):
        self.n = len(a)
        self.t = [0] * (2 * self.n)
        self.t[self.n:] = a[:]
        for i in range(self.n - 1, 0, -1):
            self.t[i] = self.t[i << 1] + self.t[i << 1 | 1]

    def update(self, pos, val):
        pos += self.n
        self.t[pos] = val
        pos >>= 1
        while pos:
            self.t[pos] = self.t[pos << 1] + self.t[pos << 1 | 1]
            pos >>= 1

    def query(self, l, r):  # inclusive [l, r]
        l += self.n
        r += self.n + 1
        res = 0
        while l < r:
            if l & 1:
                res += self.t[l]
                l += 1
            if r & 1:
                r -= 1
                res += self.t[r]
            l >>= 1
            r >>= 1
        return res
\end{lstlisting}
\end{cpblock}

\begin{cpblock}{Coordinate Compression}
\begin{lstlisting}
vals = sorted(set(a))
comp = {x: i for i, x in enumerate(vals)}
compressed = [comp[x] for x in a]
\end{lstlisting}
\end{cpblock}

\begin{cpblock}{Order Statistics with Fenwick}
Useful when values are compressed to $1..m$.
\begin{lstlisting}
def kth(bit, k):
    # returns smallest idx with prefix sum >= k
    idx = 0
    step = 1 << (bit.n.bit_length() - 1)
    while step:
        nxt = idx + step
        if nxt <= bit.n and bit.bit[nxt] < k:
            idx = nxt
            k -= bit.bit[nxt]
        step >>= 1
    return idx + 1
\end{lstlisting}
\end{cpblock}

\begin{cpblock}{Heap Queue}
\begin{lstlisting}
from heapq import heappush, heappop

pq = []
heappush(pq, (dist, node))
d, u = heappop(pq)

# max-heap trick
heappush(pq, -x)
x = -heappop(pq)
\end{lstlisting}
\end{cpblock}

\newpage

% =========================
\section{General Techniques and Tricks}
% =========================
\begin{cpblock}{Two Pointers}
\begin{lstlisting}
l = 0
for r in range(n):
    # add a[r]
    while bad(l, r):
        # remove a[l]
        l += 1
    # current window is [l, r]
\end{lstlisting}
\end{cpblock}

\begin{cpblock}{Sliding Window Sum}
\begin{lstlisting}
cur = sum(a[:k])
best = cur
for r in range(k, n):
    cur += a[r] - a[r - k]
    best = max(best, cur)
\end{lstlisting}
\end{cpblock}

\begin{cpblock}{Frequency Map}
\begin{lstlisting}
from collections import defaultdict, Counter

freq = defaultdict(int)
for x in a:
    freq[x] += 1

cnt = Counter(a)
\end{lstlisting}
\end{cpblock}

\begin{cpblock}{Sorting with Keys}
\begin{lstlisting}
points.sort()                  # lexicographic
points.sort(key=lambda p: p[1]) # by y
points.sort(key=lambda p: (p[0], -p[1]))
\end{lstlisting}
\end{cpblock}

\begin{cpblock}{Output Fast}
\begin{lstlisting}
import sys

ans = ["YES", "NO", "YES"]
sys.stdout.write("\n".join(ans))
\end{lstlisting}
\end{cpblock}

\begin{cpblock}{Memory Tricks}
\begin{itemize}
    \item Avoid huge lists of lists if a sparse structure works.
    \item For grids, store characters as strings when read-only.
    \item Use arrays of integers instead of dictionaries when indices are dense.
    \item Use \code{bytearray(n)} for boolean-like arrays when memory is tight.
\end{itemize}
\end{cpblock}


\end{document}
